\documentclass[11pt]{article}
\usepackage[T1]{fontenc}
\usepackage{lmodern}
\usepackage[letterpaper,margin=0.95in]{geometry}
\usepackage{amsmath,amssymb,amsthm,booktabs,array}
\usepackage[hidelinks]{hyperref}
\usepackage{url}
\newcommand{\R}{\mathbb R}
\newcommand{\Q}{\mathcal Q}
\newcommand{\add}{\operatorname{add}}
\newcommand{\sub}{\operatorname{sub}}
\newcommand{\mul}{\operatorname{mul}}
\newcommand{\dv}{\operatorname{div}}
\newcommand{\rint}{\operatorname{rint}}
\newcommand{\ldexp}{\operatorname{ldexp}}
\newcommand{\tp}{\operatorname{tanh}_{p}}
\newcommand{\fl}{\operatorname{fl}}
\newcommand{\code}[1]{\texttt{#1}}
\title{A Precision-Controlled Tanh Construction:\\Arithmetic Contract and Implementation Evidence}
\author{}\date{Section 3 research note 1 --- 25 September 2026}

\usepackage{microtype,fancyhdr}
\pagestyle{fancy}\fancyhf{}
\fancyhead[L]{\small True-p-bit implementation}
\fancyhead[R]{\small Section 3 research note 1}
\fancyfoot[C]{\thepage}
\setlength{\headheight}{14pt}\setlength{\emergencystretch}{2em}
\allowdisplaybreaks[2]
\begin{document}
\maketitle

\begin{abstract}
The expC11 computation constructs a fixed-geometry tanh network and fits its
readout with an explicitly rounded, $p$-significand-bit implementation of
reference LAPACK's SVD least-squares driver. The same arithmetic controls
geometry, activation evaluation, the solve, stored coefficients, and sequential
inference. This note specifies that computation, including its input boundary,
exponent range, operation order, and rank cutoff. It records the existing
implementation checks and completed chirp and Runge measurements, without
claiming a new test run or a certificate for a different recovery algorithm.
\end{abstract}

\section{What the experiment measures}

The question is whether a prescribed tanh geometry can be constructed, fitted,
and evaluated accurately when the numerical computation itself is restricted
to $p$ significant binary digits. Rounding a final coefficient vector is not
enough to answer this question: feature formation and the least-squares solve
can themselves introduce errors. ExpC11 therefore fixes a sequence of rounded
operations for all of these stages.

The objects entering the experiment are a set of target samples, a bandwidth
$\lambda$, integer dimensions, and a floating-point format. Target sampling
and bandwidth selection occur outside the controlled computation. They produce
data that are rounded before entering it. Within the computation, the centers,
features, least-squares solution, and predictions are all formed in the chosen
format. A binary64 observer measures the resulting error afterwards.

This distinction matters when connecting the experiment to Section~3's
realizability discussion. The experiment supplies an explicit implementation
and numerical evidence for one construction procedure. It is not a proof that
every solver, summation order, or parameterization has the same error, and it
does not implement the supplied reference-scaled FP64 recovery theorem. That
theorem has its own coefficient-envelope scaling and recovery/export contract.
ExpC11 uses the unscaled design matrix and the particular truncated
\code{DGELSS} procedure described below.

The earlier expC09 precision panel also has a different interpretation: its
precision restriction concerns the forward computation while the solve remains
in FP64. Its observed plateau and error constant cannot be presented as those
of an all-$p$-bit solve. The later expC13 raw-affine comparison is likewise a
separate protocol. None of its measurements are used in this note.

\section{The arithmetic contract}

\subsection{Significand precision and exponent range}
Let $\mathcal F(p,e_{\min},e_{\max})$ denote a binary format with $p$
significand bits, including the leading bit, normal exponents from $e_{\min}$
through $e_{\max}$, and gradual underflow. A positive normal value has the form
$m2^e$ with $1\le m<2$ and $p$ binary significant digits. Subnormals below
$2^{e_{\min}}$ lie on the uniform lattice with spacing
\[
 q_{\mathrm{sub}}=2^{e_{\min}-p+1}.
\]
Write $\Q$ for round-to-nearest, ties-to-even into this format. The five
arithmetic primitives are
\begin{align*}
 \add(a,b)&=\Q(a+b), & \sub(a,b)&=\Q(a-b),\\
 \mul(a,b)&=\Q(ab), & \dv(a,b)&=\Q(a/b), &
 \operatorname{sqrt}_{p}(a)&=\Q(\sqrt a).
\end{align*}
Every operand is already a format value, and each primitive rounds separately.
In particular, a multiply followed by an add is not fused. The unit roundoff
for normal results is $u_p=2^{-p}$; the primary solver cutoff is $2u_p$.

Negation, absolute value, comparison, sign transfer, and copying do not require
a rounding step. Power-of-two scaling is exact when its result is
representable; the emulator's scaling wrapper also applies format rounding at
range boundaries. The nearest-integer operation $\rint$, with ties to even,
is used only on the small values required by activation range reduction. An
integer used as a floating operand enters as $\Q(k)$. Integer loop counters,
array sizes, and indexing remain ordinary exact control information.

The precision sweep uses
\[
 \mathcal F_p=\mathcal F(p,-958,959),\qquad p=8,9,\ldots,53.
\]
This range is wide enough to separate the experiment's significand restriction
from the narrow exponent range of IEEE binary32. It also ensures that every
finite sweep-format value, including its subnormals, is exactly representable
as a normal binary64 carrier value. Native-format checks use
\[
 \mathcal F_{32}=\mathcal F(24,-126,127),\qquad
 \mathcal F_{64}=\mathcal F(53,-1022,1023).
\]
Thus the sweep point $p=24$ is not an IEEE FP32 experiment. The sweep format at
$p=53$ is also not identical to binary64 as a set of numbers. The recorded
agreement with binary64 concerns the tested computations; exponent-sensitive
behavior need not agree on arbitrary inputs.

\subsection{Constants and the input boundary}
Constants are prepared outside the model and rounded once into the format.
This includes the activation coefficients and the machine parameters used by
LAPACK. The preparation code uses exact rational arithmetic for rational
constants and high-precision evaluation with a midpoint-separation check for
the required transcendental constants. Runtime access retrieves the stored
format value; it does not recompute that constant at higher precision.

For binary64 training locations $x_i$ and supplied binary64 labels
$y_i=\fl_{64}(f(x_i))$, the model receives
\[
 \widetilde x_i=\Q(x_i),\qquad
 \widetilde y_i=\Q(y_i).
\]
It likewise receives rounded evaluation inputs $\widetilde x^e_i$, bandwidth
$\widetilde\lambda=\Q(\lambda)$, and geometry integers
$\widetilde N=\Q(N)$ and $\widetilde j=\Q(j)$. The labels are samples at
the original locations $x_i$, subsequently rounded; they are not freshly
evaluated as $f(\widetilde x_i)$ inside the model.

The bandwidth is selected offline, separately for each precision, using the
existing refined rule with tolerance $2^{1-p}$. Consequently, the experiment
does not assert that target generation or the bandwidth-selection algorithm
itself runs at $p$ bits. Its controlled boundary begins after those inputs
have been supplied and rounded.

\section{Geometry, activation, and design matrix}

\subsection{Centered features in a prescribed operation order}
Choose $N$ interior intervals and $H$ halo centers on each side. The center
indices are $j=-H,\ldots,N+H$, giving
\[
 W=N+2H+1
\]
hidden features. Geometry is formed by
\begin{equation}\label{eq:geometry}
 h=\dv(2,\widetilde N),\qquad
 c_j=\add\bigl(-1,\mul(\widetilde j,h)\bigr),\qquad
 \gamma=\dv(\widetilde\lambda,h).
\end{equation}
The feature at a format-valued input $x$ is
\begin{equation}\label{eq:feature}
 \phi_j(x)=\tp\!\left(\mul\bigl(\gamma,\sub(x,c_j)\bigr)\right).
\end{equation}
The subtraction is performed before multiplication by $\gamma$. Forming a
stored bias $b_j=\Q(-\gamma c_j)$ and evaluating
$\add(\mul(\gamma,x),b_j)$ gives a different rounding sequence, even though
the two expressions are identical over the reals. The expC11 evidence applies
to the centered form~\eqref{eq:feature}.

At small $p$, the rounding of geometry integers and centers is part of the
experiment. Distinct indices can produce equal rounded values, and a fine
training grid can have repeated rounded locations. The array still has the
prescribed number of rows and columns. One should not interpret that nominal
width as a guarantee of distinct finite-precision features or full rank.

\subsection{A tanh algorithm built from rounded primitives}
The activation is a specified algorithm $\tp$, not a call to a system
\code{tanh} routine followed by output rounding. Its construction uses
\[
 \tanh(a)=-\frac{\exp(-2a)-1}{\exp(-2a)+1},\qquad a\ge0,
\]
and evaluates the numerator as an \code{expm1}-type expression to avoid
subtracting two nearly equal numbers near the origin.

The stored constants are
\begin{align*}
 I&=\Q(1/\ln2), & S&=\Q((p+3)\ln2/2),\\
 L_h&=\Q_{\max(p-k_b,1)}(\ln2), &
 L_l&=\Q(\ln2-L_h),
\end{align*}
where $k_b$ is the bit length of $p+4$ and $\Q_q$ denotes rounding to $q$
significand bits. The shorter significand of $L_h$ supports the split range
reduction. Let $d\ge1$ be the smallest integer satisfying
\[
 \frac{0.35^d}{(d+1)!}\le2^{-(p+1)},\qquad
 a_\ell=\Q\!\left(\frac1{(\ell+1)!}\right),\quad 0\le\ell<d.
\]
These coefficients define a Horner approximation to $\exp(r)-1$ after
reduction. The degree rule is part of the specified algorithm; by itself it
is not a global floating-point error bound for the completed activation.

For input $z$, set $a=|z|$. If $a\ge S$, return $1$ or $-1$ according to
the sign of $z$. Otherwise form $u=-2a$ by exact power-of-two scaling and set
\[
 k=\rint\bigl(\mul(u,I)\bigr),\qquad
 r=\begin{cases}
 u,& k=0,\\
 \sub\bigl(\sub(u,\mul(k,L_h)),\mul(k,L_l)\bigr),& k\ne0.
 \end{cases}
\]
The small integer $k$ is represented in the working format when used in these
products. Evaluate the polynomial in the following order:
\[
 q\leftarrow a_{d-1},\qquad
 q\leftarrow\add(a_\ell,\mul(r,q))\quad
 (\ell=d-2,d-3,\ldots,0),\qquad P=\mul(r,q).
\]
Reconstruct the exponential difference and the activation by
\begin{align*}
 E&=\begin{cases}
 P,&k=0,\\
 \add\bigl(\ldexp(P,k),\sub(2^k,1)\bigr),&k\ne0,
 \end{cases}\\
 t&=-\dv(E,\add(E,2)),\qquad
 \tp(z)=\begin{cases}t,&z\ge0,\\-t,&z<0.\end{cases}
\end{align*}
Every displayed add, subtract, multiply, and divide is one rounded operation.
There is no wider accumulator inside the activation algorithm. Nor is
$\tp(z)$ claimed to equal correctly rounded mathematical $\tanh(z)$.
The existing tests compare the algorithm with an independent MPFR replay and
check its error against mathematical tanh on sampled inputs; these are
different checks.

\subsection{The sampled fitting problem}
For $M$ training points, form the $M\times n$ matrix, $n=W+1$,
\[
 A_{i,j+H+1}=\phi_j(\widetilde x_i),\qquad
 A_{i,n}=1,\qquad -H\le j\le N+H.
\]
The final column represents the output bias. The right-hand side is
$\widetilde y=(\widetilde y_1,\ldots,\widetilde y_M)^T$. Entries of $A$
are already format values when the least-squares driver receives them.
The driver performs its own rounded matrix operations; it does not hand this
matrix to a binary64 solver.

\section{The all-p-bit least-squares computation}

\subsection{Reference driver and supported branch}
The implementation ports reference LAPACK 3.12.1 \code{DGELSS}, restricted
to $M\ge n$ and one right-hand side. Floating-point expressions in the
reachable routines are mapped to the primitives above in the specified
evaluation order. Block size is fixed at one, selecting the unblocked
level-2 kernels. A blocked library implementation can change rounding order
and is not the bitwise reference for this computation.

For the tall expC11 matrix, the main stages are Householder QR, application
of its orthogonal transformations to the right-hand side, reduction of the
remaining problem to upper bidiagonal form, construction of the right
singular-vector factors, implicit-shift QR iteration on the bidiagonal matrix,
and truncated back-substitution. The port includes the scalar operations
inside these stages: norms, rotations, reflectors, matrix--vector products,
and scaling. Representative routines are \code{DGEQR2}, \code{DLARFG},
Blue's \code{DNRM2}, \code{DGEBD2}, \code{DORGL2}, \code{DBDSQR},
\code{DLARTG}, \code{DLAS2}, \code{DLASV2}, \code{DRSCL}, and
\code{DGEMV}.

Machine parameters belong to the working format. In particular,
\code{DLAMCH('E')} is $2^{-p}$ and \code{DLAMCH('P')} is $2^{1-p}$.
Safe minimum, overflow threshold, norm-scaling thresholds, and the constants
used by bidiagonal iteration are also prepared for that format. Changing only
the rounding of matrix entries while retaining binary64 machine constants
would define a different solver.

The port is deliberately not a general replacement for all \code{DGELSS}
branches. It returns an error for $M<n$, an input requiring the omitted
\code{DLASCL} rescaling branch, or an all-zero matrix. A constant bias column
rules out the all-zero design matrix here. A nonzero bidiagonal-iteration
status is propagated rather than silently accepted as a fitted model.

There is one documented arithmetic exception in solver control flow. The
reference driver's dimension crossover is computed as
\[
 \code{MNTHR}=\operatorname{INT}
       \bigl(\operatorname{REAL}(\min(M,n))\cdot1.6\bigr)
\]
in native single precision, matching reference LAPACK. It determines whether
to perform the initial QR reduction. This calculation uses matrix dimensions,
not data values; it is outside the $p$-bit numerical kernels and is explicitly
allowed by the static audit. For the reported $4801\times1025$ problem it
selects the tall-matrix QR path.

\subsection{Truncation and sensitivity cutoffs}
In exact arithmetic, a truncated SVD solve retains singular directions whose
singular values exceed a threshold, divides their transformed right-hand-side
components by those singular values, and maps back through the right singular
vectors. In expC11 each of these operations, including the decomposition, is
performed in the working arithmetic. Therefore the exact-arithmetic formula
describes the algorithmic purpose, not an identity asserting that the computed
factors are an exact SVD of $A$.

Let $\widehat\sigma_1\ge\cdots\ge\widehat\sigma_n$ be the computed
singular values, and let $\mathrm{sfmin}$ be the stored LAPACK safe minimum.
The implemented retention rule is
\begin{equation}\label{eq:threshold}
 \widehat\sigma_i>	au,\qquad
 \tau=\max\{\mul(\rho,\widehat\sigma_1),\mathrm{sfmin}\},
 \qquad \rho=2^{1-p}.
\end{equation}
Equality is discarded. This includes the safe-minimum safeguard present in
the code; the shorter description ``discard values at most
$\rho\sigma_1$'' suppresses that detail. The reported effective rank is
the number of retained computed singular values.

The chirp sensitivity study also uses $\rho=\kappa2^{-p}$ for
$\kappa\in\{1,8,32\}$. It reuses the state after \code{DBDSQR} and repeats
the cutoff-dependent scaling and final matrix--vector multiply. Since the
decomposition does not depend on \code{RCOND}, this reproduces those stages
of separate calls with the same arithmetic and decomposition. A cutoff is
not chosen afterwards to minimize the plotted primary error: the primary
curve consistently uses $\rho=2^{1-p}$.

\section{Storage, inference, and the external observer}

Write the fitted coefficients as $w_j$, $-H\le j\le N+H$, and bias $b$.
All are format values produced by the solver. For an evaluation input
$\widetilde x$, the readout starts at the bias and visits the features in
increasing center-index order:
\begin{equation}\label{eq:readout}
 s\leftarrow b,\qquad
 s\leftarrow\add\bigl(s,\mul(\phi_j(\widetilde x),w_j)\bigr),
 \quad j=-H,-H+1,\ldots,N+H.
\end{equation}
This is a sequential rounded sum. Pairwise reduction, compensated summation,
a library dot product, or an exact accumulator would change the experiment.
Saved-model evaluation reconstructs the same rounded geometry and applies
this same order.

The emulator stores values in binary64 arrays, but the storage type does not
grant additional arithmetic precision: each stored number must be exactly
representable in the chosen format, and subsequent arithmetic still goes
through its rounding primitives. Saved coefficient, prediction, and
singular-value arrays are checked for this representability. The saved
model files use binary64 containers; they are not a claim of a physically
packed $p$-bit deployment format or a measured memory saving.

With $M_e$ original evaluation locations $x^e_i$, define the reported error
\[
 E_p=\frac{\left(\sum_{i=1}^{M_e}
   |s(\Q(x^e_i))-f_{64}(x^e_i)|^2\right)^{1/2}}
 {\left(\sum_{i=1}^{M_e}|f_{64}(x^e_i)|^2\right)^{1/2}},
 \qquad f_{64}(x)=\fl_{64}(f(x)).
\]
The observer performs this measurement in binary64. In particular, it compares
against the target at the unrounded grid, so input quantization contributes
to the error. This is a finite-grid relative $\ell^2$ measurement, conventionally
called relative $L^2$ in the experiment reports. It is not a certified
continuous-domain norm, and no sampled-to-continuous theorem is invoked here.

\section{Implementation checks and their scope}

The model and solver are written once in two C headers and instantiated with
three arithmetic backends. The emulator uses a binary64 carrier to implement
the requested rounded primitives; the native variants use binary32 or
binary64 scalar types. Native builds disable contraction and automatic
vectorization, preserving the scalar operation sequence.

The emulator must avoid double rounding. A binary64 operation followed by a
naive $p$-bit rounding can choose the wrong side of a target-format midpoint.
The implementation therefore computes the carrier result together with the
sign of its residual relative to the exact result. TwoSum supplies this
information for addition; FMA residuals supply it for multiplication, division,
and square root. The residual sign resolves midpoint cases before the result
is returned in the target format. Exact power-of-two scaling protects very
small residual calculations. The internal use of FMA implements one rounded
primitive; it does not fuse a model multiply with its following model add.

The completed reports record 202 passing tests in the arithmetic and replay
suite, together with full-size native/reference comparisons. These are
historical results from those reports, not tests rerun in preparing this note.
The evidence has several distinct roles:
\begin{enumerate}
\item Primitive comparisons use MPFR for the rounded add, subtract, multiply,
divide, and square-root operations, including midpoint traps, subnormal
behavior, and tiny results. Native-format comparisons additionally check
agreement with hardware arithmetic.
\item A separate MPFR implementation, written from the specification,
replays activation constants, $\tp$, feature formation, and sequential
readout. It does not independently replay the least-squares solve. Its model
replay uses a wide exponent range without subnormalization; finite-range
behavior is covered separately by the primitive tests.
\item The solver is compared with netlib \code{SGELSS}/\code{DGELSS},
compiled from reference source with contraction disabled and block size one.
The tests cover random, rank-deficient, graded, and tanh design matrices.
The reports also record bitwise agreement on the full $4801\times1025$
problem in exact IEEE FP32 and FP64.
\item The static audit walks clang's expanded syntax tree, including every
branch, and rejects raw floating arithmetic, unrounded conversions, disallowed
math calls, or unchecked floating literals outside the rounding layer.
Its allowed findings are exactly the dimension-crossover exception described
above. Planted arithmetic and conversion leaks serve as negative controls.
\item The Runge report records representability and error-recomputation
checks for all 46 saved models, source-hash agreement with its manifest, and
bitwise saved-model replay at $p=8,24,53$.
\end{enumerate}

These checks support fidelity to a specified numerical algorithm. They do not
constitute an exhaustive proof over all inputs, a certified tanh error bound,
or a conditioning theorem. In particular, the existing tanh test samples
inputs at each precision from 8 through 53 and requires error at most three
ulps, where the ulp is measured at the true output's magnitude. The report's
stronger phrasing ``at most 2.5 ulp at every $p$'' should be read as reported
numerical evidence, not a uniform bound proved by that test.

\section{Completed measurements}

The two reported sweeps use $W=1024$, $H=24$, and $N=975$, with $4801$
training points and $8001$ evaluation points, equispaced on $[-1,1]$. The
targets are
\[
 f_{\mathrm{chirp}}(x)=\sin(8\pi(x+1)^2),\qquad
 f_{\mathrm{Runge}}(x)=\frac1{1+25x^2}.
\]
The offline selector uses frequency scales $16\pi$ and $5$, respectively.
For Runge, the search interval is $[0.03,3]$: at $p=8,9,10,11$ the selector
returns its upper limit $\lambda=3$, rather than an equality root. At
$p=12,\ldots,53$ the report records satisfaction of the stated score equation.

Table~\ref{tab:results} reproduces rounded values from the completed reports.
Both main curves use the sweep exponent range and cutoff $2^{1-p}$.
\begin{table}[ht]
\centering
\caption{Reported relative evaluation errors. The IEEE FP32 anchor has a
different exponent range from the 24-bit sweep point.}
\label{tab:results}
\begin{tabular}{lcc}
\toprule
Format & Chirp & Runge 25\\
\midrule
Sweep, $p=24$ & $6.81\times10^{-5}$ & $9.57\times10^{-6}$\\
Sweep, $p=53$ & $1.04\times10^{-13}$ & $2.50\times10^{-14}$\\
IEEE binary32 anchor & $6.39\times10^{-5}$ & not reported here\\
\bottomrule
\end{tabular}
\end{table}

For chirp, a line $E_p\approx C2^{-p}$ with slope fixed in advance was fit
by its intercept over $p=16,\ldots,40$, giving $C=1386$. Over
$p=16,\ldots,53$, the reported median of $E_p2^p$ is $1115$. The sweep
continues down to $p=53$ without the separate plateau visible in the
forward-only comparison. The latter's reported fitted constant is $19.5$,
and its FP64-solve plateau is approximately $9.3\times10^{-14}$. The data
support a one-bit-per-bit trend for this fixed-size, fixed-solver experiment;
they do not identify a universal error constant for least squares.

Low precision is a distinct regime. For chirp at $p=8,\ldots,13$, the
reported errors are $1.35,1.46,0.84,0.39,0.29,0.25$, with retained ranks
$100,232,299,575,871,943$. Larger sensitivity cutoffs reduce several spikes
in the higher-precision curve. Thus the trend must not be described as
monotone convergence at every precision or as successful accurate fitting
throughout the 8--53 range.

Runge generally has lower error through most of the sweep, including the
24-bit and 53-bit entries in the table. No convergence-rate or intercept fit
was performed for Runge. Its final two errors, $2.66\times10^{-14}$ at
$p=52$ and $2.50\times10^{-14}$ at $p=53$, do not alone establish a floor.

The chirp IEEE FP32 anchor also illustrates why the exponent range must be
reported. Its error differs by about six percent from the 24-bit sweep point;
the report records approximately $1.8\times10^8$ subnormal intermediates in
that FP32 run. This is an observed effect of the full format, not evidence
that the significand-bit specification was violated. Both source reports
retain a pending-author-review status for presentation or conclusions, even
though the cited measurement sweeps are complete.

\section{Source provenance and reproducibility boundary}

The following are actual paths relative to the repository root. They separate
the computation's specification, inspected implementation, historical test
evidence, and saved experimental results. No training run, precision sweep,
or numerical validation suite was rerun for this note.

\begin{description}\raggedright
\item[Normative contract.]
\path{experiments/expC11_true_precision_law/SPEC.md}.
The explicit safe-minimum threshold in~\eqref{eq:threshold}, zero-matrix
rejection, and dimension-crossover exception are clarified from the code.
\item[Model and solver.]
\path{src/precision/pbit_algo.h} and
\path{src/precision/lapack_gelss.h}.
\item[Arithmetic and interface.]
\path{src/precision/pbit_emul.c},
\path{src/precision/pbit_native.c}, and
\path{src/precision/pbit.py}.
\item[Independent anchors and audit.]
\path{src/precision/reference_lapack.py} and
\path{src/precision/audit.py}.
\item[Tests and model replay.]
\path{tests/test_pbit.py},
\path{tests/pbit_replay_reference.py}, and
\path{tests/test_pbit_replay_selfcheck.py}.
\item[Experiment runners.]
\path{experiments/expC11_true_precision_law/run.py} and
\path{experiments/expC11_true_precision_law/runge_comparison.py}.
\item[Chirp report.]
\path{results/checkpoint_C_geometry/expC11_true_precision_law/expC11_results.md}.
Its sibling \path{data/} directory contains \path{config.json},
\path{measurements.jsonl}, \path{summary.json}, the saved
\path{models/p*.npz}, and native/reference anchor records.
\item[Runge report.]
\path{results/checkpoint_C_geometry/expC11_true_precision_law/runge25_comparison/expC11_runge25_results.md}.
Its sibling \path{data/} directory contains the configuration and source
hashes, measurements, models, and \path{validation.json}.
\item[Distinct recovery proof supplied with the source packet.]
\path{docs/appendix_materials/2026-09-25/supplied/fp64_scaled_svd_recovery_proof.pdf}.
It is identified here to keep its numerical contract separate; the expC11
checks are not a certificate that this theorem's algorithm was implemented.
\end{description}

\end{document}
